\begin{frame}
	\frametitle{O tema comum}
	\begin{itemize}
		\item Ambas as fronteiras abandonam a \textbf{multiplicação densa em ponto flutuante} como operação fundamental --- uma pela via da \textbf{representação} (menos bits por peso), outra pela via do \textbf{tempo} (menos ativações por passo).
		\item As duas nasceram de comunidades diferentes (compressão de modelos vs. neurociência computacional) e, de forma independente, chegaram à mesma solução estrutural para treinar através de uma função não-diferenciável: o STE do BitNet e o gradiente substituto da RNP são o mesmo truque.
		\item O mesmo gargalo de fundo motiva as duas: energia dominada por \textbf{movimento de dado}, não por operações aritméticas em si --- a tabela de energia vista na abertura desta palestra.
	\end{itemize}
\end{frame}
